\SourceLeaf{114}%
\Article{32}
Die Rechenregeln für komplexe Zahlen sind allgemein bekannt: Durch Einführung der Norm wird die Division auf die Multiplikation zurückgeführt, denn es gilt
\[
\frac{a+bi}{c+di}=(a+bi)\pdot\frac{c-di}{cc+dd}=\frac{ac+bd}{cc+dd}+\frac{bc-ad}{cc+dd}\pdot i
\]

Die Quadratwurzel wird mithilfe der Formel
\[
\surd(a+bi)=\pm\left(\surd\frac{\surd(aa+bb)+a}{2}+i\surd\frac{\surd(aa+bb)-a}{2}\right)
\]
gezogen, wenn $b$ positiv ist, oder mithilfe der Formel
\[
\surd(a+bi)=\pm\left(\surd\frac{\surd(aa+bb)+a}{2}-i\surd\frac{\surd(aa+bb)-a}{2}\right)
\]
wenn $b$ negativ ist. Beim Nutzen der Umformung der komplexen Größe $a+bi$ in $r(\cos\varphi+i\sin\varphi)$ zur Erleichterung der Rechnungen brauchen wir uns hier nicht aufzuhalten.

\Article{33}
Eine komplexe ganze Zahl, die sich in zwei von den Einheiten verschiedene Faktoren\begingroup\renewcommand{\thefootnote}{*)}\footnote{Oder, was dasselbe ist, in solche, deren Normen größer als die Einheit sind.}\endgroup\ zerlegen lässt, nennen wir eine zusammengesetzte komplexe Zahl; eine Zahl, die keine solche Zerlegung in Faktoren zulässt, soll dagegen eine komplexe Primzahl heißen. Daraus folgt unmittelbar, dass jede zusammengesetzte reelle Zahl auch eine zusammengesetzte komplexe Zahl ist. Eine reelle Primzahl kann jedoch eine zusammengesetzte komplexe Zahl sein, und dies gilt für die Zahl $2$ sowie für alle positiven reellen Primzahlen der Form $4n+1$ (mit Ausnahme der Zahl $1$), da sich diese bekanntlich in zwei positive Quadrate zerlegen lassen. Zum Beispiel gilt $2=(1+i)(1-i)$, $5=(1+2i)(1-2i)$, $13=(3+2i)(3-2i)$, $17=(1+4i)(1-4i)$ usw.

Dagegen sind positive reelle Primzahlen der Form $4n+3$ stets komplexe Primzahlen. Wäre nämlich eine solche Zahl $q$ $=(a+bi)(\alpha+\beta i)$, so wäre auch $q=(a-bi)(\alpha-\beta i)$ und daher $qq=(aa+bb)(\alpha\alpha+\beta\beta)$. Nun lässt sich aber $qq$ nur auf eine einzige Weise in positive, die Einheit übersteigende Faktoren zerlegen, nämlich in $q\times q$. Daher müsste $q=aa+bb=\alpha\alpha+\beta\beta$ gelten, \mbox{Q. E. A.}; denn die Summe zweier Quadrate kann nicht die Form $4n+3$ haben.

\SourceLeaf{115}%
Die negativen reellen Zahlen behalten offenbar dieselben Benennungen wie die positiven; dasselbe gilt für die rein imaginären Zahlen.

Es bleibt somit zu zeigen, wie man unter den gemischt imaginären Zahlen die zusammengesetzten von den Primzahlen unterscheidet. Dies geschieht durch den folgenden

\textsc{Satz.}\quad\emph{Jede gemischt imaginäre ganze Zahl $a+bi$ ist eine komplexe Primzahl oder eine zusammengesetzte Zahl, je nachdem, ob ihre Norm eine reelle Primzahl oder eine zusammengesetzte Zahl ist.}

\emph{Beweis.} I. Da die Norm einer zusammengesetzten komplexen Zahl stets eine zusammengesetzte Zahl ist, muss eine komplexe Zahl, deren Norm eine reelle Primzahl ist, offenbar eine komplexe Primzahl sein. \mbox{Q. E. P.}

II. Ist aber die Norm $aa+bb$ zusammengesetzt, so sei $p$ eine positive reelle Primzahl, die sie teilt. Nun sind zwei Fälle zu unterscheiden.

1) Hat $p$ die Form $4n+3$, so kann bekanntlich $aa+bb$ nur dann durch $p$ teilbar sein, wenn $p$ beide Zahlen $a,b$ teilt. Daher ist $a+bi$ eine zusammengesetzte Zahl.

\begingroup\postdisplaypenalty=10000\relax
2) Hat $p$ nicht die Form $4n+3$, so lässt sich diese Zahl sicher in zwei Quadrate zerlegen. Wir setzen daher $p=\alpha\alpha+\beta\beta$. Da
\[
(a\alpha+b\beta)(a\alpha-b\beta)=aa(\alpha\alpha+\beta\beta)-\beta\beta(aa+bb)
\]
gilt und dieser Ausdruck somit durch $p$ teilbar ist, muss $p$ einen der Faktoren $a\alpha+b\beta$, $a\alpha-b\beta$ teilen. Da ferner
\[
(a\alpha+b\beta)^2+(b\alpha-a\beta)^2=(a\alpha-b\beta)^2+(b\alpha+a\beta)^2=(aa+bb)(\alpha\alpha+\beta\beta)
\]
gilt und dieser Ausdruck somit durch $pp$ teilbar ist, muss im ersten Fall auch $b\alpha-a\beta$, im zweiten $b\alpha+a\beta$ durch $p$ teilbar sein. Daher ist im ersten Fall
\[
\frac{a+bi}{\alpha+\beta i}=\frac{a\alpha+b\beta}{p}+\frac{b\alpha-a\beta}{p}\pdot i
\]
eine komplexe ganze Zahl, im zweiten dagegen
\[
\frac{a+bi}{\alpha-\beta i}=\frac{a\alpha-b\beta}{p}+\frac{b\alpha+a\beta}{p}\pdot i
\]
eine ganze Zahl. Da also die gegebene Zahl durch $\alpha+\beta i$ oder durch $\alpha-\beta i$ teilbar ist und die Norm des Quotienten $=\frac{aa+bb}{p}$ nach Voraussetzung von der Einheit verschieden ist, ist $a+bi$ offenbar in beiden Fällen eine zusammengesetzte komplexe Zahl. \mbox{Q. E. S.}\endgroup

\SourceLeaf{116}%
\Article{34}
Den gesamten Bereich der komplexen Primzahlen erschöpfen somit die folgenden vier Arten:

\begin{list}{}{\setlength{\leftmargin}{1.7em}\setlength{\labelwidth}{1.5em}\setlength{\labelsep}{.2em}}
\item[1)] die vier Einheiten $1,+i,-1,-i$, die wir jedoch bei der Behandlung von Primzahlen meist stillschweigend als ausgeschlossen betrachten werden;
\item[2)] die Zahl $1+i$ mit ihren drei assoziierten Zahlen $-1+i$, $-1-i$, $1-i$;
\item[3)] die positiven reellen Primzahlen der Form $4n+3$ mit jeweils drei assoziierten Zahlen;
\item[4)] die komplexen Zahlen, deren Normen reelle, die Einheit übersteigende Primzahlen der Form $4n+1$ sind. Jeder gegebenen Norm dieser Art entsprechen stets acht komplexe Primzahlen und nicht mehr, da sich eine solche Norm nur auf eine einzige Weise in zwei Quadrate zerlegen lässt.
\end{list}

\Article{35}
Wie sich die reellen ganzen Zahlen in gerade und ungerade und die ersteren wiederum in durch vier teilbare Zahlen und in das Doppelte einer ungeraden Zahl einteilen, so ergibt sich unter den komplexen Zahlen eine ebenso wesentliche Unterscheidung. Sie sind nämlich

entweder durch $1+i$ nicht teilbar, nämlich Zahlen $a+bi$, bei denen eine der beiden Zahlen $a,b$ ungerade und die andere gerade ist;

oder durch $1+i$, aber nicht durch $2$ teilbar, wenn beide Zahlen $a,b$ ungerade sind;

oder durch $2$ teilbar, wenn beide Zahlen $a,b$ gerade sind.

Die Zahlen der ersten Klasse kann man passend \emph{ungerade komplexe Zahlen}, die der zweiten \emph{halbgerade}, die der dritten \emph{gerade} nennen.

Ein Produkt aus mehreren komplexen Faktoren ist stets ungerade, wenn alle Faktoren ungerade sind; halbgerade, wenn ein Faktor halbgerade und die übrigen ungerade sind; dagegen gerade, wenn unter den Faktoren mindestens zwei halbgerade oder mindestens einer gerade ist.

Die Norm jeder ungeraden komplexen Zahl hat die Form $4n+1$; die Norm einer halbgeraden Zahl hat die Form $8n+2$; schließlich ist die Norm einer geraden Zahl das Produkt einer Zahl der Form $4n+1$ mit der Zahl $4$ oder einer höheren Potenz von zwei.

\Article{36}
Da die Beziehung zwischen vier assoziierten komplexen Zahlen der Beziehung zwischen zwei entgegengesetzten reellen Zahlen analog ist (also Zahlen, die dem absoluten Wert nach gleich und mit entgegengesetzten Vorzeichen versehen sind) und unter diesen gewöhnlich und mit Recht die positive als primär angesehen wird,\ \SourceLeaf{117}%
entsteht die Frage, ob sich unter vier assoziierten komplexen Zahlen eine ähnliche Unterscheidung festlegen lässt und ob sie als nützlich anzusehen ist. Zur Entscheidung muss man bedenken, dass das Unterscheidungsprinzip so beschaffen sein muss, dass das Produkt zweier Zahlen, die unter ihren jeweiligen assoziierten Zahlen als primär gelten, stets unter seinen eigenen assoziierten Zahlen primär ist. Doch wir gelangen bald zu der Überzeugung, dass es ein solches Prinzip überhaupt nicht gibt, wenn die Unterscheidung nicht auf ganze Zahlen beschränkt wird; eine \emph{nützliche} Unterscheidung wird sogar auf ungerade Zahlen zu beschränken sein. Für diese lässt sich das angestrebte Ziel aber auf zweierlei Weise erreichen. Nämlich:

I. Das Produkt zweier Zahlen $a+bi,a'+b'i$, bei denen $a,a'$ die Form $4n+1$ haben und $b,b'$ gerade sind, besitzt dieselbe Eigenschaft, sodass der reelle Teil $\equiv1\ (\mathrm{mod.}\ 4)$ und der imaginäre Teil gerade ist. Man sieht leicht, dass unter vier assoziierten ungeraden Zahlen nur eine diese Form hat.

\Needspace{7\baselineskip}
II. Ist eine Zahl $a+bi$ so beschaffen, dass $a-1$ und $b$ entweder beide durch vier teilbar oder beide das Doppelte einer ungeraden Zahl sind, so besitzt ihr Produkt mit einer komplexen Zahl derselben Form wieder dieselbe Form; und man sieht leicht, dass unter vier assoziierten ungeraden Zahlen nur eine diese Form hat.

Von diesen beiden annähernd gleich geeigneten Prinzipien werden wir das letztere wählen: Unter vier assoziierten ungeraden komplexen Zahlen betrachten wir diejenige als primär, die nach dem Modul $2+2i$ der positiven Einheit kongruent ist. So lassen sich mehrere bemerkenswerte Sätze kürzer aussprechen. Zum Beispiel sind die komplexen Primzahlen $-1+2i$, $-1-2i$, $+3+2i$, $+3-2i$, $+1+4i$, $+1-4i$ usw. primär, ebenso die reellen Zahlen $-3$, $-7$, $-11$, $-19$ usw., die offenbar stets mit negativem Vorzeichen zu nehmen sind. Die zu einer ungeraden primären komplexen Zahl konjugierte Zahl ist ebenfalls primär.

Für halbgerade und gerade Zahlen im Allgemeinen wäre eine ähnliche Unterscheidung zu willkürlich und zu wenig nützlich. Von den assoziierten Primzahlen $1+i$, $1-i$, $-1+i$, $-1-i$ können wir zwar eine vor den übrigen als primär auswählen; auf zusammengesetzte Zahlen werden wir eine solche Unterscheidung aber nicht ausdehnen.

\Article{37}
Sind unter den Faktoren einer zusammengesetzten komplexen Zahl solche, die selbst zusammengesetzt sind, und werden diese wiederum in ihre Faktoren zerlegt, so gelangt man offenbar schließlich zu Primfaktoren; das heißt, jede zusammengesetzte Zahl ist in Primfaktoren zerlegbar. Sind darunter solche, die nicht primär sind, so ersetzen wir jede durch\ \SourceLeaf{118}%
das Produkt der assoziierten primären Zahl mit $i,-1$ oder $-i$. So lässt sich offenbar jede zusammengesetzte komplexe Zahl $M$ auf die Form
\[
M=i^\mu A^\alpha B^\beta C^\gamma\dotsfour
\]
bringen, wobei $A,B,C$ usw. verschiedene primäre komplexe Primzahlen sind und $\mu=0,1,2$ oder $3$ ist. Für diese Zerlegung ergibt sich der Satz, dass sie nur auf eine einzige Weise möglich ist. Bei flüchtiger Betrachtung könnte dieser Satz zwar an sich selbstverständlich erscheinen, doch bedarf er jedenfalls eines Beweises. Den Weg dazu bereitet der folgende

\textsc{Satz.}\quad\emph{Das Produkt $M=A^\alpha B^\beta C^\gamma\dotsfour$, in dem $A,B,C$ usw. verschiedene primäre komplexe Primzahlen bezeichnen, kann durch keine primäre komplexe Primzahl teilbar sein, die nicht unter $A,B,C$ usw. vorkommt.}

\begingroup\postdisplaypenalty=10000\relax
\emph{Beweis.} Es sei $P$ eine primäre komplexe Primzahl, die nicht unter $A,B,C$ usw. enthalten ist, und es seien $p,a,b,c$ usw. die Normen der Zahlen $P,A,B,C$ usw. Daraus folgt leicht, dass die Norm von $M$ $=a^\alpha b^\beta c^\gamma$ usw. ist. Diese Zahl müsste daher, wenn $M$ durch $P$ teilbar wäre, durch $p$ teilbar sein. Da die einzelnen Normen entweder reelle Primzahlen (aus der Folge $2,5,13,17$ usw.) oder Quadrate reeller Primzahlen (aus der Folge $9,49,121$ usw.) sind, kann dies offenbar nur eintreten, wenn $p$ mit einer der Normen $a,b,c$ usw. identisch ist. Wir nehmen daher $p=a$ an. Da aber $P,A$ nach Voraussetzung verschiedene primäre komplexe Primzahlen sind, sieht man leicht, dass diese Bedingungen nur dann zugleich bestehen können, wenn $P,A$ konjugierte imaginäre komplexe Zahlen sind und somit $p=a$ eine ungerade reelle Primzahl ist (nicht das Quadrat einer Primzahl). Wir setzen daher $A=k+li$, $P=k-li$. Indem wir Begriff und Zeichen der Kongruenz auf komplexe ganze Zahlen ausdehnen, erhalten wir $A\equiv2k\ (\mathrm{mod.}\ P)$, woraus leicht folgt:
\[
M\equiv2^\alpha k^\alpha B^\beta C^\gamma\dotsfour\ (\mathrm{mod.}\ P)
\]
Wird also $M$ als durch $P$ teilbar angenommen, so ist auch
\[
2^\alpha k^\alpha B^\beta C^\gamma\dotsfour
\]
durch $P$ teilbar und daher die Norm dieser Zahl, die
\[
=2^{2\alpha}k^{2\alpha}b^\beta c^\gamma\dotsfour
\]
ist, durch $p$ teilbar. Da jedoch $2$ und $k$ sicher nicht durch $p$ teilbar sind, folgt daraus,\ \SourceLeaf{119}%
dass $p$ mit einer der Zahlen $b,c$ usw. identisch sein muss; sei zum Beispiel $p=b$. Daraus schließen wir aber, dass entweder $B=k+li$ oder $B=k-li$ gilt, das heißt entweder $B=A$ oder $B=P$, beides gegen die Voraussetzung.\endgroup

Aus diesem Satz folgt der andere, dass die Zerlegung in Primfaktoren nur auf eine einzige Weise möglich ist, sehr leicht durch Schlüsse, die denjenigen völlig analog sind, die wir in den \emph{Disquisitiones Arithmeticae} für reelle Zahlen verwendet haben (Artikel 16). Daher wäre es überflüssig, uns hier damit aufzuhalten.

